\documentclass[10pt,a4paper]{amsart}
\usepackage[margin=19mm]{geometry}
\usepackage{amsmath,amssymb,mathtools}
\usepackage[scr=rsfs,cal=euler]{mathalfa}
\usepackage{xfrac}
\usepackage[unicode,hidelinks,bookmarks=false]{hyperref}
\newcommand{\ie}{i.e.\ }
\newcommand{\cf}{cf.\ }
\newcommand{\resp}{respectively\ }
\newcommand{\Subsection}{\subsection}
\providecommand{\nobreakdash}{\nobreak\hbox{-}}
\setlength{\emergencystretch}{2em}
\multlinegap0pt
\usepackage{bm}
\usepackage{bbm}
\usepackage{dsfont}
\usepackage{xspace}
\usepackage{cancel}
\usepackage{mathtools}
\usepackage{amsthm}
\makeatletter
\@ifundefined{theorem}{\newtheorem{theorem}{Theorem}}{}
\@ifundefined{proposition}{\newtheorem{proposition}[theorem]{Proposition}}{}
\makeatother
\title[Jacobson identities for post-Lie algebras in positive charac]{Jacobson identities for post-Lie algebras in positive characteristic}
\author{Quentin Ehret, Nicolas Gilliers}
\date{Source: arXiv:2504.14540v2 (CC BY 4.0)}
\begin{document}
\maketitle

\noindent\emph{Source and licence.} Jacobson identities for post-Lie algebras in positive characteristic. Version arXiv:2504.14540v2 (CC BY 4.0), distributed under the Creative Commons Attribution 4.0 International licence, as recorded in the arXiv licence metadata for this exact version and shown on the abstract page https://arxiv.org/abs/2504.14540v2. Full rights notice: licenses/arxiv-2504.14540-CC-BY-4.0.txt.\par\smallskip

\noindent\textit{Editorial scope.} Counted unit: the Proposition whose source label is \texttt{prop:trucnul} in the article (source0.tex lines 755--762, source label \texttt{prop:trucnul}), delivered together with the article's own complete original proof of it (source0.tex lines 763--813, one \texttt{proof} environment of 51 lines). No mathematics is written, repaired or re-derived: statement and proof are the article's own text line for line. Required context retained uncounted: none. Every definition, display and label of those lines is retained unchanged. Omitted from this package and not redistributed: 1--754, 814--1468. Nothing inside a retained excerpt is omitted or altered: every retained body is delivered exactly as the article's lines are written, so no omission receipt of any kind accompanies this package. Citations: the retained excerpts carry no citation command at all, so no bibliography entry of the article is transcribed and no reference is redistributed. Reference adaptation: none was needed and none was made; every reference inside the delivered unit resolves inside it, so no unresolved reference or citation is produced. Printed slips kept literal: no printed word, symbol, hypothesis or number of the counted statement or of its proof was altered. Layout and class: the article's own class is \texttt{article} with packages including amsmath, amssymb, amscd, array, mathrsfs, mathtools, bm, stmaryrd, soul, titlesec, shuffle, amsthm, color, tikz-cd; the wrapper is the lane's pinned \texttt{amsart} editorial wrapper at 10pt on A4, which restates the theorem-like environments the retained text needs and adds the article's own macro definitions that the retained text needs (none), with extra wrapper packages (bm, bbm, dsfont, xspace, cancel, mathtools). The delivered numbering is the unit's own. Assets: no retained span contains a figure environment, a graphics-inclusion command, a PDF-page inclusion, a table or a TikZ picture; no image file is copied into this package, the package declares no asset dependency and no image file is redistributed with it. Licence: version arXiv:2504.14540v2 of this article is distributed under the Creative Commons Attribution 4.0 International licence, as recorded in the arXiv licence metadata for this exact version and shown on the abstract page https://arxiv.org/abs/2504.14540v2. A full rights notice is delivered at licenses/arxiv-2504.14540-CC-BY-4.0.txt. Characters: every retained excerpt is pure ASCII, so the delivered engine, XeLaTeX with its default Latin Modern text font, reports no missing character.
\begin{proposition}
\label{prop:trucnul}
Let $n\geq 1$ and $1\leq j\leq n-1$. Let $c\in {\rm Sh}(n-1,1)$ be the shuffle permutation sending $n$ to $j$. With the notations introduced before,
$$
[x_{c\cdot \ell}] = [x_{\ell_1}\cdots x_{\ell_{j-1}} x_{\ell_n} x_{\ell_j} \cdots x_{\ell_{n-1}}] = \sum_{\substack{1\leq s \leq n-1\\\sigma \in \overline{\rm UnSh}{(n-s,s)} \\ \sigma c(n)=n \\ \sigma^{-1}(n-s+1)\neq n}} (-1)^{s}[x_{\sigma c \cdot \ell}],
$$
where $\overline{\rm UnSh}{(n-s,s)}$ is the subset of unshuffle permutations of $\{1,\ldots,n-s\}\cup \{n < n-1 < \cdots < n-s+1\}$ (note the reversing in the order of the elements). 
\end{proposition}

\begin{proof} We prove the result by induction on $n\geq 1$ and $1\leq j \leq n-1$. The result trivially holds for $n=2$. Assume that the result holds for all $n\leq N$. When $j=n-1$, on one hand
\begin{align}
[x_{\ell_1}\cdots x_{\ell_n}x_{\ell_{n-1}}] = -[x_{\ell_1}\cdots x_{\ell_{n-1}}x_{\ell_n}]
\end{align}
and on the other hand, 
\begin{align}
\sum_{\substack{1\leq s \leq n-1\\ \sigma \in \overline{\rm UnSh}{(n-s,s)} \\ \sigma(n-1)=n \\ \sigma^{-1}(n-s+1)\neq n}} (-1)^{s}[x_{\sigma c\cdot \ell}] & = \sum_{\substack{2\leq s \leq n-1\\\sigma \in \overline{\rm UnSh}{(n-s,s)} \\ \sigma(n-1)=n \\ \sigma^{-1}(n-s+1)\neq n}} (-1)^{s}[x_{\sigma c\cdot \ell}] -\sum_{\substack{\sigma \in \overline{\rm UnSh}(n-1,1)\\ \sigma(n-1)=n}}[x_{\sigma c \cdot \ell}] \\ 
&= \sum_{\substack{2\leq s \leq n-1\\\sigma \in \overline{\rm UnSh}{(n-s,s)} \\ \sigma(n-1)=n \\ \sigma^{-1}(n-s+1)\neq n}} (-1)^{s}[x_{\ell_{\sigma c}}] -[x_{\ell_1},\ldots,x_{\ell_{n-1}},x_{\ell_{n}}] \label{eqn:sumtodiscussone}.
\end{align}
the sum in the right-hand side of the last equality is empty. In fact, for any $\sigma \in {\rm \overline{UnSh}}$$(n-s,s)$ with $s\geq 2$ and $\sigma(n-1)=n$, we have $\sigma^{-1}(n-1) > \sigma^{-1}(n) = n-1 $ implies $\sigma^{-1}(n-1)=n$. Hence, in the right-hand side of \eqref{eqn:sumtodiscussone}, in the partial sum over $2\leq s \leq n-1$, only the term corresponding to $s=2$ contributes. When $s=2$, the condition $\sigma^{-1}(n-2+1)\neq 2$ and $\sigma(n-1)=n$ are incompatible.
Assume the result holds for $p+1 \leq j \leq N-1$, $p\geq 1$. The Jacobi identity implies:
\begin{align*}
&[x_{\ell_1}\cdots x_{\ell_{p-1}} x_{\ell_n} x_{\ell_p} \cdots x_{\ell_{n-1}}] = [x_{\ell_1} \cdots x_{\ell_{p-1}} [x_{\ell_n},x_{\ell_p}] x_{\ell_{p+1}} \cdots x_{\ell_{n-1}}] \\ 
&\hspace{7cm}+[x_{\ell_1}\cdots x_{\ell_{p}} x_{\ell_n} x_{\ell_{p+1}} \cdots x_{\ell_{n-1}}]. 
\end{align*}
We use the induction hypothesis to infer (see the definition of $\hat{x}$ and $\hat{c}$ after the computations), 
\begin{align}
[x_{\ell_1}\cdots x_{\ell_{p-1}} [x_{\ell_n}, x_{\ell_p}] \cdots x_{\ell_{n-1}}]
&= - [x_{\ell_1}\cdots x_{\ell_{p-1}} [x_{\ell_p}, x_{\ell_n}] x_{\ell_{p+1}} \cdots x_{\ell_{n-1}}] \\
& = - \sum_{\substack{1\leq s \leq n-2\\\sigma \in \overline{\rm UnSh}{(n-1-s,s)} \\ \sigma \hat{c}(n-1)=n-1 \\ \sigma^{-1}(n-s)\neq n-1}}[\hat{x}_{(\sigma \hat{c}) \cdot \hat{\ell}_1}, \ldots, \hat{x}_{(\sigma \hat{c}) \cdot \hat{\ell}_{n-2}}, \hat{x}_{(\sigma \hat{c}) \cdot \hat{\ell}_{n-1}}] \label{eqn:proof:cancellationsun}
\end{align}
with $(\hat{x}_{\hat{\ell}_i})_{n-1 \geq i\geq 1}$ defined by $\hat{x}_{\hat{\ell}_{i}} = x_{{\ell}_{i}}$, $1\leq i \leq p-1$, $\hat{x}_{\hat{\ell}_{k}} = {x}_{{\ell}_{k+1}}$, $p \leq k \leq n-2$, $\hat{x}_{\hat{\ell}_{n-1}} = [x_{\ell_p},x_{\ell_n}]$. The permutation $\hat{c} \in \mathcal{S}_{n-1}$ is such that $\hat{x}_{\hat{c}\cdot \hat{\ell}_1}\cdots \hat{x}_{\hat{c}\cdot \hat{\ell}_{n-1}} = \hat{x}_{\hat{\ell}_1} \cdots \hat{x}_{\hat{\ell}_{p-1}}\hat{x}_{\hat{\ell}_{n-1}}\hat{x}_{\hat{\ell}_{p}} \cdots \hat{x}_{\hat{\ell}_{n-2}}.
$ We replace in \eqref{eqn:proof:cancellationsun} $\hat{x}_{\hat{\ell}_{n}}$ by its definition to obtain
\begin{align}
\eqref{eqn:proof:cancellationsun} &= - \sum_{\substack{1\leq s \leq n-2\\\hat{\sigma} \in \overline{\rm UnSh}{(n-1-s,s)} \\ \hat{\sigma} \hat{c}(n-1)=n-1 \\ \hat{\sigma}^{-1}(n-s)\neq n-1}}[\hat{x}_{\hat{\sigma} \hat{c} \cdot \hat{\ell}_1}, \ldots, \hat{x}_{\hat{\sigma} \hat{c} \cdot \hat{\ell}_{n-2}}, {x}_{{\ell}_p}, {x}_{{\ell}_n}].
\end{align}
Now, with the change of variable defined by 
$$
\hat{\sigma} \mapsto \sigma, \quad 
\begin{cases} 
\sigma(i) = \hat{\sigma}(i),& 1 \leq i \leq p-1  \\
\sigma(p)=\hat{\sigma}(p)+1 = n\\
\sigma(p+1) = n-1 \\
\sigma(j) = \hat{\sigma}(j-1) & p+2 \leq j \leq n,
\end{cases}
$$
we obtain :
\begin{align}
\eqref{eqn:proof:cancellationsun}&= - \sum_{\substack{1\leq s \leq n-1\\\sigma \in \overline{\rm UnSh}{(n-s,s)} \\ \sigma {c}(n)=n \\ \sigma c(p)=n-1 \\ \sigma^{-1}(n-s+1)\neq n}}[{x}_{(\sigma {c}) \cdot {\ell}_1}, \ldots, {x}_{(\sigma {c}) \cdot {\ell}_{n-2}}, {x}_{(\sigma c) \cdot {\ell}_{n-1}}, {x}_{(\sigma c) \cdot {\ell}_n}].
\end{align}
Moreover, we have
\begin{align}
[x_{\ell_1} \cdots x_{\ell_{p-1}}x_{\ell_p}x_{\ell_n}\cdots x_{\ell_{n-1}}] & = \sum_{\substack{1\leq s \leq n-1\\ \tilde{\sigma} \in \overline{\rm UnSh}{(n-s,s)} \\ \tilde{\sigma} \tilde{c}(n)=n \\ \tilde{\sigma}^{-1}(n-s+1)\neq n}} [x_{(\sigma \tilde{c})\cdot\ell_1},\ldots,x_{(\sigma \tilde{c})\cdot\ell_n}] \\
& = \sum_{\substack{1\leq s \leq n-1\\\sigma \in \overline{\rm UnSh}{(n-s,s)} \\ \tilde{\sigma}\tau{c}(n)=n \\ \tilde{\sigma}^{-1}(n-s+1)\neq n}} [x_{(\sigma\tau {c})\cdot\ell_1},\ldots,x_{(\sigma \tau {c})\cdot\ell_n}].\label{eqn:bidule}
\end{align}
Next, we make the change of variables $\tilde{\sigma} \mapsto \sigma,\,  \sigma = \tilde{\sigma}\tau$. Since $\sigma^{-1}(n-k) = \tau \tilde{\sigma}^{-1}(n-k) = \tilde{\sigma}^{-1}(n-k)$ for $0 < k\leq s$ as $\tilde{\sigma}^{-1}(n-k) >  p+1$ and $\sigma^{-1}(n) = p$, we get $\sigma \in \overline{\rm UnSh}{(n-s,s)}$. Also, $\sigma c (p) = \sigma(p+1) = \tilde{\sigma}(p)\leq n-s$. We infer
\begin{align}
\eqref{eqn:bidule}& = \sum_{\substack{1\leq s \leq n-1\\\sigma \in \overline{\rm UnSh}{(n-s,s)} \\ \sigma{c}(n)=n \\ \sigma c(p)\leq n-s \\ \sigma^{-1}(n-s+1)\neq n}} [x_{(\sigma {c})\cdot\ell_1},\ldots,x_{(\sigma {c})\cdot\ell_n}].
\end{align}
The induction step is complete and the proof is done.
\end{proof}

\end{document}
